% TOP_PROBLEM: {"id": 1382, "title": "Orchard-Planting Problem", "category_id": 6, "category": {"id": 6, "name": "geometry", "display_name": "Geometry", "description": "Euclidean and non-Euclidean geometry, geometric structures.", "slug": "geometry", "order_index": 6, "created_at": "2026-07-31T15:26:25.671Z"}, "status": "open"}
% FIRST_POSTED: null
% ENTRY_KIND: statement_only
% Imported from: batch06.tex; UnsolvedMath ID 1382
\documentclass[11pt]{article}
\usepackage[margin=25mm]{geometry}
\usepackage{fontspec}
\setmainfont{Latin Modern Roman}
\usepackage{amsmath,amssymb,mathtools}
\usepackage{unicode-math}
\setmathfont{Latin Modern Math}
\usepackage{xurl,hyperref}
\hypersetup{colorlinks=true,urlcolor=blue}
\setcounter{secnumdepth}{0}
\setlength{\parindent}{0pt}
\setlength{\parskip}{5pt}
\setlength{\emergencystretch}{3em}
\newcommand{\sref}[2]{\hyperlink{source:#1:#2}{[S#2]}}
\newcommand{\cataloguescope}[1]{\paragraph{Recorded target.}#1}
\begin{document}
% UnsolvedMath ID 1382; source catalogue code GEOM-023
\setcounter{section}{1381}
\section{Orchard-Planting Problem}\label{1382}
{\small\sffamily UnsolvedMath ID 1382 · Geometry}
\subsection{Definitions and mathematical statement}

\paragraph{Shared notation.}$\mathbb N=\{1,2,\ldots\}$, and $\mathbb Z,\mathbb Q,\mathbb R,\mathbb C$ denote the integers, rationals, reals and complex numbers. Any other conventions are specified locally.


For a finite set $P$ of distinct points in the Euclidean plane, let
\[
 t_3(P)=\#\{\ell:\ \ell\text{ is a straight line and }|P\cap\ell|=3\}.
\]
Only lines containing exactly three selected points are counted. No general-position assumption or prohibition on lines containing four or more selected points is imposed.
\paragraph{Statement recorded in the supplied source.}
For every integer $n\ge1$, determine
\[
 T_3(n)=\max\{t_3(P):P\subset\mathbb R^2,\ |P|=n\}.
\]
\sref{1382}{2}
\subsection{Short English statement}
What is the largest number of straight rows containing exactly three trees obtainable by planting n trees at distinct planar points?
\subsection{Sources}
\begin{description}
\item[\hypertarget{source:1382:1}{S1}] ULAM.AI and the UnsolvedMath Contributors, \emph{UnsolvedMath: A Curated Collection of Open Mathematics Problems}, record 1382 (GEOM-023). CC BY 4.0. \url{https://huggingface.co/datasets/ulamai/UnsolvedMath}.
\item[\hypertarget{source:1382:2}{S2}] Ben Green and Terence Tao, \emph{On sets defining few ordinary lines} (2013), arXiv:1208.4714v3, Theorem 1.3 and Section 9 (the orchard problem). \url{https://arxiv.org/abs/1208.4714v3}.
\end{description}
\label{end:1382}
\end{document}
